\subsection{Depth Peeling}\index{depthpeeling}
\paragraph{}Ez az \ac{OIT} technika~\cite{Everitt01interactiveorder-independent} a képteret több ``rétegre'' osztja fel, majd ezeket a tárolt $\alpha$-érték szerint összekeveri a kimeneti kép elkészítéséhez.
\paragraph{}Az első réteg eltárolja a nézőpontból látható $RGB\alpha$ értékeket, a második mindazt, amit az első kitakart, és így tovább~\autoref{fig:depthpeeling}. Ily módon ``implicit'' sorrendezéssel pontos eredményt állít elő $n$ rétegből~\ref{eq:simpleblend}. A többmenetes rajzolás költsége összetettebb jelenetek esetén eltörpül a sorbarendezés (idő-)költsége melett.
\paragraph{}Az eljárás teljesítménye javítható különféle módokon, például:
\begin{itemize}
	\label{list:depthpeelvariants}
	\item Az átlátszatlan objektumok mélységének egy külön előmenetbeli\index{prepass, előmenet} felrajzolása egy textúrába\index{textúra}, amelyet felhasználhatunk a képtér hátulról előre haladó hámozásakor, így nem kell az algoritmus futtatni olyan fragmenseken\index{fragment, fragmens} amelyek a nézőpont elől el vannak takarva.
	\item Az $n$ réteg külön-külön tárolása helyett létrehozunk egy gyüjtő- és egy pásztázó framebuffer-t\index{framebuffer}. A pásztázó buffer tartalma közvetlen a képrajzolás után belekeverhető a gyüjtőbe. Ez a megoldás jól kombinálható az előzővel, ugyanis ehhez is hátulról előre kell bejárni a képteret.
\end{itemize}

\vfill
\begin{figure}[bp]
	\centering
	\begin{subfigure}{\textwidth}
		\centering
		\begin{tikzpicture}[rectangle, minimum size = 20pt, node distance=2cm]
			\node[] (0) {
				\includegraphics[scale=0.05]{camera-equipment-image-svgrepo-com.eps}
			};
			\node[draw, right of = 0, label=below:1.réteg] (1) {$RGB\alpha$};
			\draw[->] (0) to node[above] {} (1);
			\node[right of = 1] (2) {\ldots};
			\draw[->] (1) to node[above] {} (2);
			\node[draw, right of = 2, label=below:$n$. réteg] (3) {$RGB\alpha$};
			\draw[->] (2) to node[above] {} (3);
			\label{fig:depthpeeling}
		\end{tikzpicture}
		\caption{A képteret rajzolással sorbarendezzük, sorbarendezéses rajzolás helyett.}
	\end{subfigure}
	\vfill%
	\begin{subfigure}{\textwidth}
		\begin{equation}
			\label{eq:simpleblend}
			\Large \rgbdddot{C_{AB}} = \rgbdddot{C_A} + (1 - \alpha_A) \cdot \rgbdddot{C_B}
		\end{equation}
		\vspace*{-1.5em}\caption{A keverés maga egyszerű.}
	\end{subfigure}
\end{figure}